% ============================================================
% sec/03stage2.tex - Etapa 2 (04/10 - 06/11)
% ============================================================
\section{Etapa 2: modelado, explicabilidad y evaluaci\'on}
\label{sec:stage2}

\notebox[Objetivo de la etapa]{Ejecutar el dise\~no de la Etapa 1: construir
la tabla de estados, entrenar la escalera de modelos B0 a B4 bajo un
protocolo sin fuga, evaluar una sola vez en prueba, explicar el modelo final
con SHAP, analizar la calibraci\'on (m\'odulo del grupo de tres) y escribir
la metodolog\'ia definitiva y los resultados del art\'iculo. Peso: 45\,\%.}

\subsection{Resumen de ramas}

\begin{table}[H]
  \centering
  \caption{Ramas de la Etapa 2 y su resultado.}
  \label{tab:stage2summary}
  \footnotesize
  \begin{tabularx}{\textwidth}{@{}P{5.4cm} P{1.7cm} P{2.3cm} Y@{}}
    \toprule
    \textbf{Rama} & \textbf{Fechas} & \textbf{Responsables} & \textbf{Resultado}\\
    \midrule
    {\scriptsize\code{03DataCollectionAndEnvironmentSetup}} & 04/10--11/10 & Adriel & 55\,196 partidas en Parquet, 62 caracter\'isticas, particiones A/B/C y 24 pruebas de fuga\\
    {\scriptsize\code{04MainAIModelTrainingAndDevelopment}} & 11/10--30/10 & Daniel y Adriel & B0 a B4 entrenados, LightGBM, ablaciones, H2 y modelo final exportado\\
    \code{05FinalModelChecks} & 30/10--05/11 & Sebasti\'an y Adriel & Protocolo congelado, evaluaci\'on \'unica en prueba, calibraci\'on, SHAP y veredicto de H1 a H3\\
    \code{06PendientesDelArticulo} & 05/11--06/11 & Adriel & Art\'iculo corregido, con resultados, discusi\'on y conclusiones\\
    \bottomrule
  \end{tabularx}
\end{table}

\subsection{C\'omo se construy\'o el sistema, parte por parte}

La Fig.~\ref{fig:flujo} muestra el orden en que se construy\'o cada pieza.
Cada bloque se cerr\'o y se verific\'o antes de empezar el siguiente: no se
entren\'o ning\'un modelo hasta tener las particiones probadas, y no se abri\'o
el conjunto de prueba hasta tener todos los modelos, la calibraci\'on y la
revisi\'on de fuga listos.

\begin{figure}[H]
  \centering
  \begin{tikzpicture}[
    font=\footnotesize,
    caja/.style={draw=tecsteel, fill=boxbg, rounded corners=2pt,
                 text width=3.0cm, align=center, minimum height=1.0cm},
    rama/.style={font=\scriptsize\itshape, text=labelgray},
    flecha/.style={-{Stealth[length=2mm]}, thick, tecsteel}]
    \node[caja] (c) {Corpus p\'ublico\\10 d\'ias de JSON};
    \node[caja, right=0.6cm of c] (t) {Tabla de estados\\62 caracter\'isticas};
    \node[caja, right=0.6cm of t] (p) {Particiones A, B, C\\por partida};
    \node[caja, below=1.0cm of p] (b) {B0 $\rightarrow$ B1 $\rightarrow$ B2\\$\rightarrow$ B3 $\rightarrow$ B4};
    \node[caja, left=0.6cm of b] (x) {LightGBM, ablaciones\\H2, modelo final};
    \node[caja, left=0.6cm of x] (e) {Evaluaci\'on final\\una sola vez en prueba};
    \draw[flecha] (c) -- (t);
    \draw[flecha] (t) -- (p);
    \draw[flecha] (p) -- (b);
    \draw[flecha] (b) -- (x);
    \draw[flecha] (x) -- (e);
    \node[rama, above=0.05cm of t] {rama 03};
    \node[rama, below=0.05cm of x] {rama 04};
    \node[rama, below=0.05cm of e] {rama 05};
  \end{tikzpicture}
  \caption{Orden de construcci\'on de la Etapa 2.}
  \label{fig:flujo}
\end{figure}

% ============================================================
\subsection{Rama 03: datos y entorno}
\label{subsec:rama03}

\jentry{2026-10-04}{Apertura de la etapa y lectura de una partida real}%
       {Adriel}{03DataCollectionAndEnvironmentSetup}

Se tradujo la documentaci\'on al espa\~nol, se reemplaz\'o el plan \'agil de
la Etapa 1 por \code{PlanDeDesarrollo.md} y se prepar\'o el entorno de
Colab. Antes de escribir el extractor se inspeccion\'o el JSON de una
partida completa para entender el formato real:

\begin{itemize}
  \item En cada paso hay dos vistas. El jugador \code{ACTIVE} debe decidir;
        el \code{INACTIVE} espera y su vista va un paso atrasada.
  \item La mano del oponente llega como \code{null} y los premios como una
        lista de \code{null}: la informaci\'on oculta ya viene protegida.
  \item Como todos los premios son nulos, los premios restantes se cuentan
        por el \emph{largo} de la lista, no por los elementos no nulos.
\end{itemize}

\decision{Escribir la l\'ogica en \code{src/\allowbreak{}ptcg\_\allowbreak{}wp/\allowbreak{}*.py} y dejar los
cuadernos solo para ejecutar y mostrar resultados.}{El c\'odigo en m\'odulos
se puede probar con \code{pytest}, revisar en un PR e importar desde el
agente de la Etapa 3.}{L\'ogica dentro de las celdas: imposible de probar y
de reutilizar.}

\milestone{Primer extractor: 61 caracter\'isticas propio/oponente, marca de
preparaci\'on y descarte de empates.}{Commit ``feat(features): extraer
estados del jugador activo y a\~nadirlos a la ingesta''.}

% ------------------------------------------------------------

\jentry{2026-10-11}{Estados resueltos, pruebas de fuga y almac\'en en Kaggle}%
       {Adriel}{03DataCollectionAndEnvironmentSetup}

Al revisar las partidas completas se encontr\'o que algunos estados ya
muestran el desenlace: un jugador sin premios por tomar o sin Pok\'emon en
juego. Se agregaron la columna \code{decided} y el filtro \code{usable}, el
n\'umero de turno p\'ublico (\code{game\_\allowbreak{}turn}) y los identificadores de los
Pok\'emon activos. Los agentes se seudonimizaron con SHA-256 truncado.

\decision{Excluir los estados de preparaci\'on y los ya resueltos de toda
tabla de modelado.}{Un estado resuelto tiene el desenlace a la vista;
incluirlo inflar\'ia el desempe\~no con casos triviales. Son el 0.2\,\% de los
estados.}{Conservarlos: el modelo aprender\'ia a leer un tablero ya
terminado.}

\decision{Probar la ausencia de fuga con pruebas autom\'aticas que
\emph{intentan} provocarla.}{Las pruebas llenan la mano rival y los premios
con cartas reales, cambian el resultado y las recompensas, y cortan los
pasos futuros. En todos los casos las caracter\'isticas deben quedar
id\'enticas. Si alguien cambia el extractor y filtra algo oculto, la prueba
falla.}{Revisar a mano el c\'odigo del extractor: no protege contra cambios
futuros.}

\decision{Guardar las tablas procesadas en un dataset privado de Kaggle,
nunca en Git.}{Son gigabytes de datos y el repositorio debe guardar la
receta, no los datos. Kaggle permite que los tres integrantes bajen
exactamente las mismas tablas.}{Versionar los Parquet en Git o que cada uno
reconstruyera las tablas (unos 12 minutos por d\'ia).}

\blocker{Los compa\~neros recib\'ian un error 404 al bajar el dataset
completo, y la herramienta \code{kaggle} informaba \'exito aunque la
publicaci\'on fallara.}{Nadie m\'as que Adriel pod\'ia trabajar con las
tablas.}{Se pas\'o a descargar archivo por archivo y \code{store.py} empez\'o
a leer la salida de \code{kaggle} y fallar si no se bajaba nada.}

\aiuse{Claude}{implementar la ingesta, el extractor, las pruebas y el
almac\'en en Kaggle, y diagnosticar los errores de Colab}{Adriel ejecut\'o
cada celda y compar\'o los conteos con lo esperado antes de cada commit}

% ------------------------------------------------------------

\jentry{2026-10-11}{Particiones, cobertura de cartas y diccionario de datos}%
       {Adriel}{03DataCollectionAndEnvironmentSetup}

Se construyeron las particiones de los tres escenarios, agrupadas por
partida, con la semilla 2026.

\decision{Escenario~B por d\'ias completos: del 18 al 25 de junio a
entrenamiento, el 7 de julio a validaci\'on, el 11 a calibraci\'on, y el 15 y
el 22 a prueba. El 17 de junio queda fuera.}{Partir dentro de un d\'ia
mezclar\'ia el presente con el futuro. El 17 de junio fue un arranque parcial
de la competencia. Ning\'un d\'ia mostr\'o una concentraci\'on an\'omala de
mazos (m\'inimo 274 mazos distintos frente a una mediana cercana a 450).}{
Cortes por fecha y hora dentro de un d\'ia; reservar un d\'ia an\'omalo, que
al final no existi\'o.}

\decision{Escenario~C por arquetipos: toda partida con alguno de 33
Pok\'emon principales reservados va a prueba (15.1\,\% de las partidas).}{El
corpus no trae mazos completos. El Pok\'emon que m\'as tiempo pasa activo
identifica bien el arquetipo y permite reservar estrategias enteras.}{
Reservar mazos completos: casi no se repiten, as\'i que la reserva ser\'ia
arbitraria.}

\decision{Cobertura de cartas: en A y C toda carta activa debe aparecer en
entrenamiento; si no, se mueve a entrenamiento la partida m\'as corta que la
contiene. En B no se mueve nada.}{Sin esta regla, la ablaci\'on de H2 se
evaluar\'ia sobre cartas que el modelo nunca vio. En B, mover partidas
futuras a entrenamiento ser\'ia fuga temporal, as\'i que las cartas nuevas
(0.86 a 1.20\,\% de las filas) se aceptan como parte del cambio que B
mide.}{Regenerar la partici\'on con otra semilla hasta que la cobertura
saliera completa.}

Al revisar la tabla se encontr\'o que la tasa de victoria de la posici\'on 0
sub\'ia de 0.507 a 0.559 con los d\'ias. La auditor\'ia mostr\'o la causa: el
jugador que empieza ocupa la posici\'on 0 en el 76.4\,\% de las partidas del
17 de junio y en el 100\,\% de las del 22 de julio. No es una ventaja de la
posici\'on sino de mover primero, y por eso el modelo recibe
\code{first\_\allowbreak{}player\_\allowbreak{}is\_\allowbreak{}me} y no la posici\'on.

\milestone{Tabla final verificada: 55\,196 partidas, 7\,540\,389 estados,
7\,286\,386 filas utilizables, 119 empates excluidos, 0 partidas an\'omalas y
24 pruebas automáticas.}{PR \#8, \code{DiccionarioDeDatos.md} y cuaderno
\code{01\_\allowbreak{}exploracion}.}

% ============================================================
\subsection{Rama 04: entrenamiento de la escalera de modelos}
\label{subsec:rama04}

\jentry{2026-10-12}{Acuerdos t\'ecnicos antes de entrenar}%
       {Daniel}{04MainAIModelTrainingAndDevelopment}

Antes de escribir una sola l\'inea de entrenamiento, Daniel fij\'o por
escrito las reglas de modelado (\code{12-\allowbreak{}10-\allowbreak{}26\_\allowbreak{}AcuerdosTecnicosModelado.md}).

\decision{Cada partida pesa lo mismo, en el entrenamiento y en las
m\'etricas: cada estado recibe peso $1/n$, con $n$ los estados de su
partida.}{Sin esa correcci\'on, una partida larga pesar\'ia diez veces m\'as
que una corta solo por durar m\'as.}{Pesar cada fila por igual.}

\decision{B1 es una heur\'istica fija:
$(\textit{premios rival} - \textit{premios propios}) +
\textit{diferencia de PV}/\textit{unidad de PV} +
\textit{diferencia de banca}$. Solo la pendiente y el intercepto que la
convierten en probabilidad se ajustan con entrenamiento.}{B1 debe responder
si hace falta aprendizaje autom\'atico; si sus pesos se aprendieran, ser\'ia
otro modelo supervisado.}{Ajustar los pesos de la heur\'istica con las
etiquetas.}

\decision{Seleccionar con validaci\'on cruzada agrupada por partida
(\code{GroupKFold}), comparar finalistas en validaci\'on y dejar prueba
cerrada.}{Cada conjunto tiene un solo uso: entrenamiento ajusta, validaci\'on
elige, calibraci\'on recalibra y prueba reporta.}{\code{KFold} o
\code{StratifiedKFold} sin grupos, que mezclar\'ian estados de una partida
entre pliegues.}

\decision{La parada temprana de B4 vigila un grupo de partidas completas,
separado dentro del entrenamiento. Se fija \code{scikit-learn==1.7.2}, la
primera versi\'on que acepta datos de vigilancia externos.}{La vigilancia
interna por defecto separa filas al azar y pondr\'ia estados de una misma
partida a ambos lados.}{La vigilancia por defecto de scikit-learn.}

\aiuse{Claude}{discutir las alternativas de protocolo}{las reglas quedaron
escritas y fechadas antes de ejecutar}

% ------------------------------------------------------------

\jentry{2026-10-12 a 2026-10-13}{B0 y B1}%
       {Daniel}{04MainAIModelTrainingAndDevelopment}

\blocker{Cargar todas las columnas de todos los d\'ias agotaba la RAM de
Colab.}{El cuaderno de B0 y B1 no terminaba.}{Leer solo las columnas
necesarias y recorrer las tablas d\'ia por d\'ia.}

\milestone{B0 y B1 medidos en validaci\'on para A, B y C. AUC de B1: 0.677,
0.668 y 0.671.}{Cuaderno \code{02\_\allowbreak{}baselines\_\allowbreak{}B0\_\allowbreak{}B1}.}

% ------------------------------------------------------------

\jentry{2026-10-15 a 2026-10-18}{B2 con todo el entrenamiento en la RAM de Colab}%
       {Daniel}{04MainAIModelTrainingAndDevelopment}

La primera versi\'on de B2 cargaba la matriz completa y fallaba por memoria.
Una muestra de 5\,000 partidas lo hac\'ia caber, pero desperdiciaba la mayor
parte de los datos.

\decision{Entrenar B2 por lotes: \code{StandardScaler.partial\_\allowbreak{}fit} con
pesos $1/n$ y \code{SGDClassifier} con p\'erdida log\'istica y
penalizaci\'on L2, leyendo los Parquet en grupos.}{Permite usar todo el
entrenamiento sin cargarlo en memoria. Se convirti\'o $C$ en
$\alpha = 1/(C\cdot G)$, con $G$ las partidas, para conservar el mismo
objetivo ponderado de la regresi\'on log\'istica.}{Entrenar con una muestra:
el baseline no habr\'ia visto los mismos datos que los dem\'as modelos.}

\blocker{Las pruebas de equivalencia detectaron que, con
\code{average=True}, scikit-learn dejaba de actualizar el intercepto al
encadenar \code{partial\_\allowbreak{}fit}.}{B2 habr\'ia dado un modelo distinto del que
dice describir.}{Se promediaron los coeficientes a mano al terminar cada
\'epoca. Una medici\'on sint\'etica de 2 millones de estados confirm\'o un pico
de solo 279~MiB.}

% ------------------------------------------------------------

\jentry{2026-10-19}{Primera b\'usqueda de B2 y lectura de la falta de convergencia}%
       {Daniel}{04MainAIModelTrainingAndDevelopment}

Se compararon nueve valores de $C$, de $10^{-4}$ a $10^{4}$, con cinco
pliegues por partida en A, B y C. Desde $C = 0.1$ las m\'etricas se aplanan.
Ning\'un pliegue cumpli\'o el criterio de convergencia (gradiente
$\le 10^{-3}$).

\decision{No relajar la tolerancia para que el indicador dijera
``convergi\'o''. Volver a entrenar solo las finalistas ($C = 0.1$ y
$C = 1$) con hasta 600 \'epocas.}{``0 de 5'' significa que el criterio no
pudo certificar el m\'inimo, no que el modelo haya fallado. Cambiar el umbral
solo para cambiar el indicador habr\'ia ocultado el problema.}{Bajar la
tolerancia o reportar la b\'usqueda como concluyente.}

% ------------------------------------------------------------

\jentry{2026-10-19 a 2026-10-23}{Reanudaci\'on y auditor\'ia del costo}%
       {Daniel}{04MainAIModelTrainingAndDevelopment}

Colab cortaba las sesiones a mitad de los entrenamientos. Se agregaron
puntos de control por pliegue y por \textit{trial}, y un orquestador que
ejecuta B2, B3 y B4 en secuencia. Cuando B3 empez\'o a tardar horas sin
avanzar, se audit\'o el c\'odigo (\code{23-\allowbreak{}10-\allowbreak{}26\_\allowbreak{}AuditoriaBaselinesYCoste.md})
y se encontraron cuatro problemas:

\begin{itemize}
  \item Al reanudar, el muestreador de Optuna se recreaba con la misma
        semilla y \textbf{repet\'ia configuraciones ya evaluadas}. Se
        reprodujo: el tercer \textit{trial} repiti\'o los par\'ametros del
        primero.
  \item Con 10 \textit{trials} y 10 de arranque aleatorio, \textbf{la
        b\'usqueda bayesiana nunca llegaba a guiarse} por los resultados.
  \item Tres estudios por m\'etrica multiplicaban el costo: 900 ajustes de
        B3 y 450 de B4, sin contar los finales.
  \item La base SQLite de Optuna viv\'ia directamente en Google Drive, con
        escrituras cada minuto.
\end{itemize}

\decision{Versi\'on \code{efficient\_\allowbreak{}v2}: cinco \textit{trials} por
estudio, tres de arranque aleatorio, espacio de b\'usqueda acotado (por
ejemplo, bosques de 25 a 100 \'arboles con profundidad 8 o 12), SQLite en
disco local con respaldos consistentes y reanudaci\'on que reencola el
\textit{trial} interrumpido.}{Elimina las configuraciones de costo extremo
que hac\'ian inviable el presupuesto, sin perder los resultados ya
obtenidos. Las rutas nuevas no se mezclan con las anteriores.}{Seguir con el
presupuesto original, que no terminaba en semanas.}

% ------------------------------------------------------------

\jentry{2026-10-25}{Selecci\'on final de B2}%
       {Daniel}{04MainAIModelTrainingAndDevelopment}

\begin{table}[H]
  \centering
  \caption{B2: finalistas en validaci\'on (\code{25-\allowbreak{}10-\allowbreak{}26\_\allowbreak{}B2SeleccionFinalC.md}).}
  \label{tab:b2final}
  \footnotesize
  \begin{tabular}{@{}lcccc@{}}
    \toprule
    \textbf{Escenario} & $\mathbf{C}$ & \textbf{AUC} & \textbf{Log loss} & \textbf{Brier}\\
    \midrule
    A & \textbf{0.1} & \textbf{0.7607} & \textbf{0.5774} & \textbf{0.1975}\\
    A & 1.0 & 0.7607 & 0.5775 & 0.1975\\
    B & \textbf{0.1} & \textbf{0.7281} & \textbf{0.6309} & \textbf{0.2140}\\
    B & 1.0 & 0.7278 & 0.6324 & 0.2143\\
    C & \textbf{0.1} & \textbf{0.7603} & \textbf{0.5786} & \textbf{0.1977}\\
    C & 1.0 & 0.7602 & 0.5787 & 0.1977\\
    \bottomrule
  \end{tabular}
\end{table}

\decision{$C = 0.1$ en los tres escenarios, y cerrar B2 con el presupuesto
invertido.}{Gana en las tres m\'etricas en los tres escenarios. Los
gradientes finales (0.0035 a 0.0072) no certifican el m\'inimo, y eso se
reporta junto a las m\'etricas.}{Seguir aumentando \'epocas sin l\'imite, o
usar prueba para resolver el empate.}

% ------------------------------------------------------------

\jentry{2026-10-25 a 2026-10-26}{B4: una sola b\'usqueda, sin fuga entre escenarios}%
       {Daniel y Adriel}{04MainAIModelTrainingAndDevelopment}

Daniel prepar\'o un cuaderno independiente para B4. Adriel revis\'o su
l\'ogica antes de lanzar el entrenamiento completo y encontr\'o un riesgo:
buscar hiperpar\'ametros por separado en cada escenario usaba partidas que
eran validaci\'on o prueba en \emph{otro} escenario.

\decision{Una sola b\'usqueda (Optuna TPE, 20 \textit{trials}, tres
pliegues por partida) sobre las 11\,285 partidas que son entrenamiento en A,
B y C a la vez. Un solo estudio optimiza la p\'erdida logar\'itmica y
registra AUC y Brier.}{As\'i ning\'un hiperpar\'ametro vio partidas de
validaci\'on, calibraci\'on o prueba de ning\'un escenario. Un estudio en
lugar de tres divide el costo entre tres.}{Una b\'usqueda por escenario, que
filtraba informaci\'on entre escenarios.}

\blocker{En el primer reentrenamiento final, con tope de 2000 \'arboles, los
tres escenarios llegaron al tope: el modelo segu\'ia mejorando.}{El n\'umero
de \'arboles no lo decid\'ia la parada temprana sino el tope.}{Se subi\'o el
tope a 4000. La parada temprana se detuvo sola en 2888, 2370 y 2714
\'arboles, y las m\'etricas cambiaron menos de 0.002. Se us\'o un servidor
con 8 n\'ucleos y 51~GiB de RAM.}

\milestone{B4 en validaci\'on: AUC 0.835 (A), 0.786 (B) y 0.836 (C).}{Cuaderno
\code{05\_\allowbreak{}entrenamiento\_\allowbreak{}B4\_\allowbreak{}independiente\_\allowbreak{}original.ipynb}, conservado como
evidencia.}

% ------------------------------------------------------------

\jentry{2026-10-28}{Selecci\'on de B3, definici\'on de H2 y protocolo de LightGBM}%
       {Daniel}{04MainAIModelTrainingAndDevelopment}

\decision{Usar en B3 una sola configuraci\'on com\'un, la de mayor
macro-AUC de validaci\'on: el promedio simple del AUC de A, B y C.}{Se
necesitaba una configuraci\'on compartida, y cada escenario debe pesar lo
mismo. El macro-AUC solo elige; los resultados siguen report\'andose por
escenario.}{Elegir una configuraci\'on distinta por escenario, o promediar
log loss y Brier en un puntaje arbitrario.}

\begin{table}[H]
  \centering
  \caption{B3: finalistas con los mismos hiperpar\'ametros en los tres
  escenarios (AUC en validaci\'on).}
  \label{tab:b3sel}
  \footnotesize
  \begin{tabular}{@{}lcccc@{}}
    \toprule
    \textbf{Finalista} & \textbf{A} & \textbf{B} & \textbf{C} & \textbf{Macro-AUC}\\
    \midrule
    \'Arbol / log loss & 0.7740 & 0.7343 & 0.7791 & 0.7624\\
    \'Arbol / Brier & 0.7767 & 0.7344 & 0.7810 & 0.7641\\
    Bosque / Brier & 0.8016 & 0.7621 & 0.8038 & 0.7892\\
    \textbf{Bosque / log loss} & \textbf{0.8029} & \textbf{0.7609} & \textbf{0.8063} & \textbf{0.7901}\\
    \bottomrule
  \end{tabular}
\end{table}

\decision{Definir H2 como la comparaci\'on entre B4 con las 62
caracter\'isticas y B4 restringido a cuatro entradas: turno, qui\'en empez\'o
y los identificadores de los dos Pok\'emon activos.}{La hip\'otesis pregunta
si la se\~nal temprana se explica solo por el orden de turno y la identidad
del mazo. Los activos son la mejor aproximaci\'on a esa identidad disponible
en cada estado sin usar datos del resto de la partida.}{Usar los mazos
completos, que solo se conocen al terminar la partida.}

\decision{Cambiar B4 por LightGBM solo si, en cada escenario, el intervalo
del 95\,\% de la diferencia de AUC queda entero por encima de cero y la
p\'erdida logar\'itmica y el Brier no empeoran. Presupuesto igual al de B4.}{
La regla se fij\'o antes de ejecutar, para que no pudiera ajustarse al
resultado.}{Elegir el que saliera mejor sin regla previa.}

\aiuse{Claude}{proponer la regla de LightGBM (decisi\'on delegada)}{Daniel la
document\'o antes de ejecutar}

% ------------------------------------------------------------

\jentry{2026-10-28 a 2026-10-29}{Auditor\'ia de B0 a B3 y Kaggle como almac\'en \'unico}%
       {Daniel y Adriel}{04MainAIModelTrainingAndDevelopment}

Se audit\'o que B0 a B3 cumplieran el contrato de modelado del diccionario
de datos (\code{29-\allowbreak{}10-\allowbreak{}26\_\allowbreak{}AuditoriaBaselinesB0B3.md}): carga con filas
utilizables, transformaciones ajustadas solo con entrenamiento, validaci\'on
agrupada, roles de cada conjunto, 62 caracter\'isticas sin identificadores y
m\'etricas por etapa. Todos cumplieron. Tambi\'en se unific\'o la
evaluaci\'on en una sola funci\'on para todos los modelos.

\blocker{Los resultados viv\'ian en Google Drive personal y en archivos
\code{.zip} descargados a mano, y los \code{.joblib} no llegaban a GitHub por
el \code{.gitignore}.}{Cada integrante ve\'ia una versi\'on distinta de los
resultados.}{Adriel reemplaz\'o Drive por el dataset privado de Kaggle como
almac\'en \'unico (\code{store.py}). Cada archivo publicado registra su
SHA-256, el commit que lo produjo y la fecha, y cada publicaci\'on conserva
todo lo anterior.}

\decision{Cuadernos sin interruptores: cada celda declara los archivos que
produce y llama a \code{store.asegurar}. Si existen en Kaggle, los baja y
verifica; si faltan, entrena y publica.}{Un interruptor olvidado en
\code{True} puede reentrenar y sobrescribir horas de trabajo; uno en
\code{False}, saltarse un paso. Sin interruptores, ejecutar la celda siempre
hace lo correcto.}{Variables como \code{EJECUTAR = False} o
\code{CONSTRUIR = True}.}

% ------------------------------------------------------------

\jentry{2026-10-29 a 2026-10-30}{LightGBM, ablaciones, H2 y modelo final}%
       {Adriel}{04MainAIModelTrainingAndDevelopment}

\milestone{LightGBM no cumple la regla y se mantiene B4: AUC LightGBM/B4 de
0.8346/0.8346 (A), 0.7849/0.7863 (B) y 0.8352/0.8355 (C). Dos
implementaciones independientes coinciden, as\'i que el resultado no depende
de la librer\'ia.}{Carpeta \code{lightgbm/\allowbreak{}} en Kaggle.}

\milestone{Ablaciones: sin recursos y tablero el AUC cae de 0.05 a 0.06; sin
progreso, como mucho 0.003; sin diferenciales no cae.}{Carpeta
\code{b4/\allowbreak{}ablaciones/\allowbreak{}}.}

\decision{No quitar los diferenciales de B4 aunque no aporten.}{Elegir
caracter\'isticas mirando validaci\'on ser\'ia una selecci\'on sobre el
mismo conjunto que se usa para comparar.}{Simplificar el modelo con esas
cifras.}

\blocker{\code{HistGradientBoosting} admite como m\'aximo 255 categor\'ias por
columna, y H2 necesita 359 cartas activas.}{Las salidas r\'apidas, agrupar o
descartar las cartas poco frecuentes, perd\'ian informaci\'on de identidad,
que es justo lo que H2 mide.}{Adriel exigi\'o conservar las 359 cartas. Cada identificador se
reparti\'o sin p\'erdida en bloques de hasta 255 categor\'ias, ordenados por
frecuencia. Esto no oblig\'o a reentrenar B4.}

\decision{El modelo final para la Etapa 3 es B4 entrenado en A, expuesto en
\code{ptcg\_\allowbreak{}wp.predictor.PredictorVictoria}.}{A tiene la poblaci\'on m\'as
amplia y reciente; B y C sirven para medir generalizaci\'on. El predictor
recibe la observaci\'on del juego y produce las mismas caracter\'isticas, en
el mismo orden, que la tabla de entrenamiento.}{Elegir el modelo de B o C, o
entrenar uno nuevo con todos los datos.}

\milestone{Escalera completa en validaci\'on (Tabla~\ref{tab:escaleraval}) y
modelo exportado que reproduce las m\'etricas de validaci\'on de B4 con
diferencia m\'axima de $5.6\times10^{-16}$.}{PR \#9 y carpeta
\code{modelo\_\allowbreak{}final/\allowbreak{}} en Kaggle.}

\begin{table}[H]
  \centering
  \caption{AUC en validaci\'on de la escalera de modelos al cerrar la rama 04.}
  \label{tab:escaleraval}
  \footnotesize
  \begin{tabular}{@{}lccccc@{}}
    \toprule
    \textbf{Escenario} & \textbf{B0} & \textbf{B1} & \textbf{B2} & \textbf{B3} & \textbf{B4}\\
    \midrule
    A & 0.500 & 0.677 & 0.761 & 0.803 & 0.835\\
    B & 0.500 & 0.668 & 0.728 & 0.761 & 0.786\\
    C & 0.500 & 0.671 & 0.760 & 0.806 & 0.836\\
    \bottomrule
  \end{tabular}
\end{table}

\aiuse{Claude}{revisar la b\'usqueda de B4, implementar el almac\'en en
Kaggle, la codificaci\'on por bloques de H2 y el predictor}{Adriel prob\'o
cada cambio en Colab y rechaz\'o las propuestas con interruptores, con
dependencia de Drive o con p\'erdida de categor\'ias}

% ============================================================
\subsection{Rama 05: evaluaci\'on final}
\label{subsec:rama05}

\jentry{2026-11-02}{Diagn\'ostico y protocolo congelado antes de abrir prueba}%
       {Sebasti\'an}{05FinalModelChecks}

Sebasti\'an reconstruy\'o el estado del proyecto
(\code{02-\allowbreak{}11-\allowbreak{}26\_\allowbreak{}DiagnosticoFinalModelChecks.md}) y encontr\'o tres riesgos
de m\'etodo antes de escribir el protocolo:

\begin{itemize}
  \item El orden sugerido del cuaderno dejaba SHAP y la revisi\'on de fuga
        para despu\'es de prueba. As\'i no pod\'ian funcionar como barrera.
  \item Comparar directamente la ca\'ida del AUC con la del log loss
        favorec\'ia a H3, porque est\'an en escalas distintas.
  \item Comparar etapas con todas las partidas mezcla duraci\'on y progreso:
        las partidas cortas nunca llegan a la etapa final.
\end{itemize}

\decision{Congelar el protocolo con un commit antes de abrir prueba
(\code{02-\allowbreak{}11-\allowbreak{}26\_\allowbreak{}ProtocoloEvaluacionFinal.md}), con este orden: verificar
la carga de modelos, SHAP y revisi\'on de fuga, decidir la calibraci\'on solo
con calibraci\'on, curvas de aprendizaje y, solo entonces, prueba.}{Si una
regla se escribe despu\'es de ver el resultado, puede acomodarse a \'el. La
celda de prueba se niega a correr si falta alguna barrera.}{Evaluar prueba
primero y decidir despu\'es.}

\decision{Decidir H3 con la descomposici\'on de la p\'erdida en calibraci\'on
(MCB) y ordenamiento (DSC), sobre las mismas 1\,405 partidas de prueba de A
y B.}{Deja ambas p\'erdidas en las mismas unidades. Usar las mismas
partidas a\'isla el efecto del entrenamiento del de la composici\'on del
conjunto.}{Comparar porcentajes de ca\'ida de AUC y de log loss.}

\decision{Controlar la duraci\'on en H1 con una cohorte: solo las partidas
que llegan a las tres etapas.}{Si el AUC sube de inicio a final dentro de las
mismas partidas, el aumento no viene de que las partidas largas sean m\'as
f\'aciles.}{Comparar etapas con todas las partidas.}

\decision{Revisi\'on de fuga con SHAP: aprueba si ninguna columna prohibida
participa, si ninguna caracter\'istica concentra m\'as del 50\,\% de la
importancia y si ninguna caracter\'istica sola alcanza un AUC de 0.85.}{Una
variable que domina el modelo suele ser la firma de una fuga. Los umbrales se
fijaron antes de mirar.}{Revisar las gr\'aficas SHAP a ojo.}

\aiuse{Claude Code}{diagn\'ostico, protocolo, implementaci\'on del cuaderno 06
y sus pruebas sint\'eticas}{Sebasti\'an revis\'o y aprob\'o el protocolo; la
herramienta no tuvo acceso a los datos}

% ------------------------------------------------------------

\jentry{2026-11-02 a 2026-11-03}{Ejecuci\'on del cuaderno 06}%
       {Adriel}{05FinalModelChecks}

Adriel ejecut\'o el cuaderno de arriba abajo en Colab Pro, sin que nadie m\'as
publicara en Kaggle al mismo tiempo. Antes de prueba pasaron todas las
barreras: 109 pruebas autom\'aticas, reproducci\'on exacta de las m\'etricas de
validaci\'on (diferencias $\le 4.4\times10^{-16}$ en log loss y Brier, y
$\le 4.8\times10^{-7}$ en AUC) y revisi\'on de fuga aprobada (importancia
m\'axima 0.113; mejor caracter\'istica sola con AUC 0.690).

\blocker{La publicaci\'on del an\'alisis de errores se marc\'o como fallida
aunque hab\'ia funcionado: los archivos \code{errores\_\allowbreak{}*.csv} contienen la
palabra ``error'', y \code{store.py} la confund\'ia con un mensaje de
Kaggle.}{La celda se detuvo despu\'es de abrir prueba.}{Se corrigi\'o la
detecci\'on para leer solo los mensajes de error reales y se agregaron
pruebas. Al volver a ejecutar, lo publicado se baj\'o sin recalcular: prueba
no se volvi\'o a predecir.}

\blocker{Justo despu\'es de publicar, el \'indice del dataset a\'un no mostraba
la versi\'on nueva a otros servidores.}{Una celda pod\'ia creer que un
archivo faltaba.}{El almac\'en guarda una marca de publicaci\'on y prefiere la
copia m\'as reciente; si Kaggle no responde, se espera unos minutos y se
reintenta.}

% ------------------------------------------------------------

\jentry{2026-11-04 a 2026-11-05}{Resultados y veredicto}%
       {Sebasti\'an y Adriel}{05FinalModelChecks}

\begin{table}[H]
  \centering
  \caption{AUC de B4 en prueba (intervalo del 95\,\%) y veredicto de las hip\'otesis.}
  \label{tab:resumenprueba}
  \footnotesize
  \begin{tabularx}{\textwidth}{@{}l c c c Y@{}}
    \toprule
    & \textbf{A} & \textbf{B} & \textbf{C} & \textbf{Lectura}\\
    \midrule
    AUC global & 0.833 & 0.789 & 0.787 & Supera a B0--B3 con intervalos que excluyen el cero\\
    AUC inicio / final & 0.684 / 0.937 & 0.630 / 0.915 & 0.623 / 0.919 & La predicci\'on mejora con la partida\\
    H1 & Sostenida & Sostenida & Sostenida & Cohorte: +0.32 de inicio a final\\
    H2 & Refutada & Sostenida & Sostenida & El estado aporta cuando cambia la poblaci\'on\\
    H3 & \multicolumn{3}{c}{Refutada} & El ordenamiento se degrada m\'as que la calibraci\'on\\
    \bottomrule
  \end{tabularx}
\end{table}

\begin{figure}[H]
  \centering
  \begin{subfigure}{0.48\textwidth}
    \centering
    \safeimage[width=\linewidth]{../../02Papers/00MainPaper/img/res_auc_por_turno.png}
    \caption{AUC de B4 por turno (A).}
  \end{subfigure}\hfill
  \begin{subfigure}{0.48\textwidth}
    \centering
    \safeimage[width=\linewidth]{../../02Papers/00MainPaper/img/res_fiabilidad.png}
    \caption{Diagramas de fiabilidad de B4.}
  \end{subfigure}
  \caption{Dos resultados de la evaluaci\'on final, tomados del cuaderno 06.}
  \label{fig:resultados}
\end{figure}

\decision{Usar el modelo final sin calibrador.}{En A, B4 ya est\'a calibrado
(ECE 0.002) y la isot\'onica lo empeora un poco. En B la isot\'onica baja el
ECE de 0.037 a 0.018, pero la misma estrategia empeor\'o a B3: corregir bajo
deriva no es fiable en general. La elecci\'on se hizo con calibraci\'on, antes
de abrir prueba.}{Calibrar siempre por defecto.}

\decision{No agregar una comprobaci\'on extra de procedencia de los
archivos de prueba.}{El cuaderno 06 se versiona con sus salidas, y el
\'indice del dataset ya registra el SHA-256, el commit y la fecha de cada
archivo. La comprobaci\'on no a\~nad\'ia informaci\'on nueva.}{Una celda
adicional que repitiera esa verificaci\'on.}

\milestone{Evaluaci\'on final cerrada: 51 archivos de resultados en
\code{evaluacion\_\allowbreak{}final/\allowbreak{}}, verificados con SHA-256.}{PR \#10,
\code{04-\allowbreak{}11-\allowbreak{}26\_\allowbreak{}ResultadosEvaluacionFinal.md} y cuaderno \code{06} con
salidas.}

% ============================================================
\subsection{Rama 06: art\'iculo y cierre}
\label{subsec:rama06}

\jentry{2026-11-05 a 2026-11-06}{Correcci\'on del art\'iculo y documentaci\'on final}%
       {Adriel}{06PendientesDelArticulo}

Se corrigieron en el art\'iculo los datos que la Etapa 2 cambi\'o: el corpus
final (55\,196 partidas), el conteo de premios, los d\'ias del escenario~B, la
definici\'on del escenario~C por arquetipos, la cobertura de cartas medida,
el sesgo de la posici\'on 0 y la URL real del conjunto de datos. Despu\'es se
escribieron las secciones de resultados, discusi\'on y conclusiones, y se
actualiz\'o la documentaci\'on del repositorio.

\decision{Declarar en el art\'iculo las desviaciones respecto del dise\~no:
una sola semilla, ablaciones sobre tres familias, sin coeficientes de B2 ni
importancia por permutaci\'on, y escenario~C por arquetipos.}{Un cambio no
declarado har\'ia que el lector crea que se ejecut\'o el dise\~no
original.}{Omitirlas o reescribir la metodolog\'ia como si siempre hubiera
sido as\'i.}

\milestone{Art\'iculo con metodolog\'ia definitiva, resultados, discusi\'on y
conclusiones.}{Commits de la rama \code{06PendientesDelArticulo}.}

% ============================================================
\subsection{C\'omo protegimos los datos}
\label{subsec:proteccion}

La Tabla~\ref{tab:proteccion} re\'une las medidas contra la fuga de
informaci\'on. Cada una responde a una v\'ia concreta por la que el modelo
podr\'ia ver algo que no deber\'ia.

\begin{table}[H]
  \centering
  \caption{V\'ias de fuga y c\'omo se cerr\'o cada una.}
  \label{tab:proteccion}
  \footnotesize
  \begin{tabularx}{\textwidth}{@{}P{4.2cm} Y P{3.0cm}@{}}
    \toprule
    \textbf{V\'ia de fuga} & \textbf{Medida} & \textbf{Comprobaci\'on}\\
    \midrule
    Estados de una partida a ambos lados de una partici\'on & Toda partici\'on, pliegue, vigilancia e intervalo se agrupa por partida & Pruebas de particiones\\
    Informaci\'on oculta (mano rival, premios) & Solo se lee la vista del jugador que act\'ua; los campos ocultos nunca entran & Pruebas que llenan lo oculto y exigen X id\'entico\\
    Resultado, recompensas y pasos futuros & Columnas prohibidas por c\'odigo & Pruebas de fuga\\
    Estados con el desenlace a la vista & Filtro \code{usable} (sin preparaci\'on ni estados resueltos) & Prueba de carga\\
    Futuro en el escenario B & D\'ias completos en orden; no se mueven partidas futuras a entrenamiento & Prueba de cronolog\'ia\\
    Hiperpar\'ametros que ven otros escenarios & B\'usqueda de B4 solo con las 11\,285 partidas de entrenamiento comunes & Prueba de la b\'usqueda\\
    Preprocesamiento ajustado con todos los datos & Escalado y cortes de etapa ajustados solo con entrenamiento & Auditor\'ia de B0--B3\\
    Decisiones tomadas con prueba & Protocolo congelado; prueba se abre una vez, despu\'es de todas las barreras & Celda de prueba con barreras\\
    Una variable que filtra el resultado & Revisi\'on SHAP con umbrales previos & Aprobada antes de prueba\\
    \bottomrule
  \end{tabularx}
\end{table}

\subsection{C\'omo aseguramos que todo funcionara}
\label{subsec:verificacion}

\begin{itemize}
  \item \textbf{Pruebas autom\'aticas.} Pasaron de 24 al cerrar la rama de
        datos a 109 al ejecutar la evaluaci\'on final. Cubren fuga,
        particiones, evaluaci\'on, almac\'en, entrenamiento y predictor, con
        datos sint\'eticos.
  \item \textbf{Valor esperado en cada celda.} Cada celda anota lo que debe
        mostrar (por ejemplo, ``55\,196 partidas''). Si la salida no
        coincid\'ia, no se hac\'ia commit.
  \item \textbf{Reproducci\'on antes de avanzar.} El modelo exportado y el
        cuaderno 06 deb\'ian reproducir las m\'etricas ya publicadas antes de
        usarse.
  \item \textbf{Trazabilidad de cada archivo.} Cada resultado en Kaggle
        guarda su SHA-256, el commit que lo produjo y la fecha.
  \item \textbf{Reanudaci\'on segura.} Puntos de control por pliegue,
        \textit{trial} y escenario, con manifiestos que impiden mezclar
        configuraciones distintas.
  \item \textbf{Una sola persona publica.} Al ejecutar algo nuevo, se
        avisaba al grupo para que nadie m\'as publicara al mismo tiempo.
\end{itemize}

% ============================================================
\subsection{Registro de experimentos}
\label{subsec:registro}

Semilla 2026 en todas las corridas, salvo B3 (semilla efectiva 2030). Las
cifras de AUC van en el orden A / B / C.

\begin{table}[H]
  \centering
  \caption{Corridas de la escalera de modelos (B0 a B4).}
  \label{tab:stage2runs}
  \footnotesize
  \begin{tabularx}{\textwidth}{@{}P{3.0cm} P{4.6cm} P{1.5cm} Y@{}}
    \toprule
    \textbf{Corrida} & \textbf{Configuraci\'on} & \textbf{Fecha} & \textbf{Resultado}\\
    \midrule
    B0 y B1 & Prevalencia; heur\'istica de recursos con pendiente e intercepto ajustados en entrenamiento & 13/10 & AUC de B1 en validaci\'on 0.668--0.677\\
    B2, b\'usqueda & 9 valores de $C$, 5 pliegues, todo el entrenamiento & 19/10 & M\'etricas planas desde $C = 0.1$; 0 de 5 pliegues convergen\\
    B2, finalistas & $C \in \{0.1, 1\}$, hasta 600 \'epocas & 25/10 & $C = 0.1$; AUC 0.761 / 0.728 / 0.760\\
    B3, \code{efficient\_\allowbreak{}v2} & \'Arbol y bosque, 5 \textit{trials} por estudio & 23--28/10 & Bosque com\'un por macro-AUC; AUC 0.803 / 0.761 / 0.806\\
    B4, b\'usqueda & Optuna TPE, 20 \textit{trials}, 3 pliegues, 11\,285 partidas & 26/10 & Configuraci\'on del \textit{trial} 20\\
    B4, tope 2000 & Reentrenamiento final & 26/10 & Los tres escenarios llegaron al tope\\
    B4, tope 4000 & Reentrenamiento final & 26/10 & 2888 / 2370 / 2714 \'arboles; AUC 0.835 / 0.786 / 0.836\\
    \bottomrule
  \end{tabularx}
\end{table}

\begin{table}[H]
  \centering
  \caption{Corridas de experimentos y evaluaci\'on final.}
  \label{tab:stage2runs2}
  \footnotesize
  \begin{tabularx}{\textwidth}{@{}P{3.0cm} P{4.6cm} P{1.5cm} Y@{}}
    \toprule
    \textbf{Corrida} & \textbf{Configuraci\'on} & \textbf{Fecha} & \textbf{Resultado}\\
    \midrule
    LightGBM & Mismo presupuesto que B4 & 29--30/10 & No cumple la regla; se mantiene B4\\
    Ablaciones & B4 sin cada familia & 29--30/10 & Recursos y tablero es la familia esencial\\
    H2 & B4 restringido a 4 entradas, 359 cartas sin p\'erdida & 30/10 & AUC 0.733 / 0.614 / 0.743 en validaci\'on\\
    Curvas de aprendizaje & 25, 50, 75 y 100\,\% del entrenamiento & 03/11 & El \'ultimo 25\,\% suma $\le$ 0.004\\
    Evaluaci\'on final & Una sola vez en prueba & 03/11 & AUC 0.833 / 0.789 / 0.787\\
    \bottomrule
  \end{tabularx}
\end{table}


% ============================================================
\subsection{Decisiones consolidadas}

\begin{table}[H]
  \centering
  \caption{Decisiones de la Etapa 2: datos y baselines (11/10 a 25/10).}
  \label{tab:stage2decisions}
  \footnotesize
  \begin{tabularx}{\textwidth}{@{}l P{4.2cm} Y P{2.4cm}@{}}
    \toprule
    \textbf{Fecha} & \textbf{Decisi\'on} & \textbf{Por qu\'e} & \textbf{Art\'iculo}\\
    \midrule
    11/10 & Excluir preparaci\'on y estados resueltos & Tienen el desenlace a la vista & Metodolog\'ia\\
    11/10 & Escenario B por d\'ias completos, sin el 17 de junio & Evita mezclar presente y futuro & Metodolog\'ia\\
    11/10 & Escenario C por 33 arquetipos & El corpus no trae mazos completos & Metodolog\'ia\\
    11/10 & Cobertura de cartas en A y C, no en B & En B ser\'ia fuga temporal & Metodolog\'ia\\
    12/10 & Peso $1/n$ por partida & Una partida larga no debe pesar m\'as & Metodolog\'ia\\
    12/10 & B1 con pesos fijos & Responde si hace falta aprendizaje & Resultados\\
    18/10 & B2 por lotes con todo el entrenamiento & Cabe en Colab sin muestrear & Resultados\\
    19/10 & No relajar la tolerancia de B2 & No ocultar la falta de convergencia & Resultados\\
    23/10 & B\'usquedas \code{efficient\_\allowbreak{}v2} & El presupuesto original no terminaba & Resultados\\
    25/10 & $C = 0.1$ en B2 & Mejor en las tres m\'etricas & Resultados\\
    \bottomrule
  \end{tabularx}
\end{table}

\begin{table}[H]
  \centering
  \caption{Decisiones de la Etapa 2: B4, experimentos y evaluaci\'on final (26/10 a 06/11).}
  \label{tab:stage2decisions2}
  \footnotesize
  \begin{tabularx}{\textwidth}{@{}l P{4.2cm} Y P{2.4cm}@{}}
    \toprule
    \textbf{Fecha} & \textbf{Decisi\'on} & \textbf{Por qu\'e} & \textbf{Art\'iculo}\\
    \midrule
    26/10 & Una sola b\'usqueda de B4 en partidas comunes & Sin fuga entre escenarios & Resultados\\
    26/10 & Tope de 4000 \'arboles & El tope de 2000 se alcanzaba & Resultados\\
    28/10 & Macro-AUC para B3 & Una configuraci\'on com\'un, escenarios con igual peso & Resultados\\
    28/10 & Regla de LightGBM previa & No acomodar la regla al resultado & Resultados\\
    29/10 & Kaggle como almac\'en \'unico y cuadernos sin interruptores & Una sola versi\'on de los resultados, con trazabilidad & Resultados\\
    30/10 & H2 con las 359 cartas sin p\'erdida & H2 mide justo la identidad & Resultados\\
    30/10 & Modelo final: B4 de A & Poblaci\'on m\'as amplia & Conclusiones\\
    02/11 & Protocolo congelado y barreras antes de prueba & Prueba se usa una vez & Resultados\\
    02/11 & H3 con descomposici\'on MCB/DSC & Escalas comparables & Resultados\\
    04/11 & Modelo final sin calibrador & Decidido con calibraci\'on, no con prueba & Resultados\\
    06/11 & Declarar las desviaciones del dise\~no & Transparencia & Resultados\\
    \bottomrule
  \end{tabularx}
\end{table}


\subsection{Retrospectiva}

\begin{table}[H]
  \centering
  \caption{Retrospectiva de la Etapa 2.}
  \label{tab:stage2retro}
  \footnotesize
  \begin{tabularx}{\textwidth}{@{}P{4.0cm} Y@{}}
    \toprule
    \textbf{Pregunta} & \textbf{Respuesta}\\
    \midrule
    ¿Qu\'e funcion\'o? & Cerrar cada bloque antes del siguiente. Escribir las reglas antes de ejecutar (acuerdos, protocolo de LightGBM, protocolo final). Probar la fuga con pruebas que intentan provocarla. Kaggle como almac\'en \'unico con trazabilidad por archivo.\\
    ¿Qu\'e nos bloque\'o? & La RAM y los cortes de sesi\'on de Colab, y el costo de las b\'usquedas de B3 y B4 con el presupuesto original. El almac\'en en Drive, que dejaba versiones distintas en cada integrante.\\
    ¿Qu\'e estimamos mal? & El tiempo de entrenamiento: B2 y B3 tomaron unos cuatro y tres d\'ias de calendario por las reanudaciones. Tambi\'en el tope de \'arboles de B4, que hubo que duplicar.\\
    ¿Qu\'e cambia para la Etapa 3? & Partir del predictor ya exportado y probado, mantener el almac\'en en Kaggle y las celdas sin interruptores, y fijar por escrito c\'omo se verificar\'a el consumo del modelo antes de construir el agente.\\
    \bottomrule
  \end{tabularx}
\end{table}